\documentclass[12pt,a4paper,titlepage]{article}

\usepackage{xcolor}
\usepackage{titlesec}
\usepackage{amsmath}
\usepackage[portuguese]{babel}
\usepackage[colorlinks=true,allcolors=teal]{hyperref}

\newcommand{\sectionbreak}{\clearpage}

\title{Equações de 1ª Ordem Exatas}
\author{Fernando de Lima Gonçalves}
\date{Belo Horizonte\\22/10/2026}

\begin{document}
\maketitle
\tableofcontents

\section{Introdução}
Considere uma equação diferencia de 1ª ordem escrita na forma geral

\[
    M(x,y) + N(x,y) = \frac{dy}{dx}
\]

Seja $\Psi(x,y) = C$ uma equação que descreve $y$ explicitamente como uma função de $x$, isto é,

\[
    \Psi(x,y(x)) = C
\]

é uma equação nas variáveis $x$ e $y$ em que $y$ é uma função implícita de $x$.
Aplicando a derivação implícita sobre a equação obtemos:

\[
    \frac{d}{dx}(\Psi(x,y)) = 0
\]

Nosso objetivo é:

\[
    \frac{\partial \Psi}{\partial x} = M(x,y) \quad , \quad \frac{\partial \Psi}{\partial y} = N(x,y)
\]

Nesse caso devemos ter:

\[
    \frac{\partial M}{\partial y} = \frac{\partial N}{\partial x}
\]

\section{Teorema}
Sejam $M,N$ duas funções de duas variáveis contínuas.
A E.D.O de 1ª ordem 

\[
    M(x,y) + N(x,y) y' = 0
\]

é exata se, e somente se, 

\[
    \frac{\partial M}{\partial y} = \frac{\partial N}{\partial x}
\]

Nesse caso, existe uma função contínua $\Psi(x,y)$ tal que 

\[
    \frac{\partial \Psi}{\partial x} = M(x,y) \quad \text{e, } \quad \frac{\partial \Psi}{\partial y} = N(x,y)
\]

A equação \fbox{$\Psi(x,y) = C$} descreve $y$ implicitamente como função de $x$ e é solução da E.D.O

\section{Lista de Equações}
\[
    \boxed{\Psi(x,y(x)) = C}
\]

\[
    \boxed{\frac{\partial \Psi}{\partial x} = M(x,y)} \quad , \quad \boxed{\frac{\partial \Psi}{\partial y} = N(x,y)}
\]

\[
    \boxed{\frac{\partial \Psi}{\partial x} + \frac{\partial \Psi}{\partial y} \frac{dy}{dx} = 0}
\]

\[
    \boxed{\frac{d}{dx}(\Psi(x,y)) = 0}
\]

\section{Exercícios}
\begin{enumerate}
    \item Identifique se as equações a seguir são \textbf{exatas}. Em caso afirmativo, resolva-as
    \begin{enumerate}
        \item $(2xy + 3) + (x^2-1)y' = 0$
        \item $(ye^{xy}) + 2x)dx + (xe^{xy} - 2y)dy = 0$
        \item $(ycos(x) + 2xe^y) + (sen(x) + x^2e^y - 1)y' = 0$
    \end{enumerate}
\end{enumerate}
\end{document}
